will be forced to work with the Chow ring, since it is on that ring that the homomorphism @@M0@@ is naturally defined as a homomorphism with values in @@M1@@; compare (XIV 4.1). Unfortunately, no theory of specialization for the Chow ring has yet been developed (cf. 7.14.) that would allow one to write commutative squares analogous to (7.13.9) and (7.13.10) in terms of Chow rings and cohomology with coefficients @@M2@@. All one can say when @@M3@@ is smooth and proper is that there is a commutative square
\sgaStableTarget{sga6.e10.eq.7-13-12}@@M4@@
analogous to (7.13.9), and a corresponding square analogous to (7.13.10), which we leave to the reader to write down, where @@M5@@ denotes the group of algebraic cycles (modulo no relation at all) of codimension @@M6@@, where the horizontal arrows are the ``associated cohomology class'' homomorphisms defined in (SGA 5 IV), and where @@M7@@ denotes the specialization homomorphism for algebraic cycles, associating with a cycle @@M8@@ the cycle @@M9@@, where @@M10@@ is the schematic closure of @@M11@@ in @@M12@@ (cf. 7.14). This commutativity follows from the following fact, established in (SGA 5 IV): if @@M13@@ is a closed subscheme of @@M14@@ of codimension @@M15@@ on every fiber, one can associate with it a cohomology class @@M16@@ whose formation commutes with every base change and, in particular, such that its restriction to @@M17@@ (resp. @@M18@@) equals @@M19@@ (resp. @@M20@@). Of course, when one knows that the homomorphism @@M21@@ in (7.13.12) carries cycles rationally equivalent to zero to cycles rationally equivalent to zero (@@M22@@ being projective over @@M23@@),
